\documentclass{article}
\usepackage{amsmath,hyperref}
\title{When two repairs need the same fixture}
\date{}
\begin{document}
\maketitle
Assume two new duties must each be mounted on the same central
fixture, then moved to their own private fixture. The latter transfer
holds both fixtures until physical completion. This extends the
two-tree vacancy model by allowing simultaneous pending repair
requests and a late change in the authority to vacate. Neither a
completed mount nor permission to transfer clears the central
fixture. The next repair can begin only after actual transfer.
In the model a real local arbiter serializes access to the center.
The leaf capacity and safety interlock are assumed, not derived.
Uncontrolled losses, preemption and cross-leaf demands are excluded.

\paragraph{Known permissions.}
Duty $i$ can begin mounting at $r_i$, takes $h_i>0$ ticks to mount,
may depart no sooner than its effective authorization time $a_i$,
and takes $c_i>0$ further ticks to physically release the center.
Each leaf is dedicated and available. With the center free at $t$,
earliest release is
\[
 F_i(t)=\max\{\max(t,r_i)+h_i,a_i\}+c_i
       =\max\{t,q_i\}+p_i,quad
 q_i=\max\{r_i,a_i-h_i\},\quad p_i=h_i+c_i .
\]
The difference in $q_i$ is truncated at zero on a natural clock.
For a fixed sequence these monotone release functions compose;
an existing in-flight mount contributes its actual residual
release time to the tail, not a fresh job. An independently
fixed partial order or separate job deadlines are not modeled.
To minimize final release time, sort by nondecreasing $q_i$.
For $q_i\le q_j$, starting $i$ before $j$ gives
$F_j(F_i(t))\le F_i(F_j(t))$. This follows by expanding both
maxima and then gives the sorting result by adjacent swaps.
It is the ordinary single-machine makespan problem with release
dates, with an effective release date derived from the mechanism.
A fixed deadline suffices for both duties only if their ordering
achieves the joint bound, not if each independently meets it.

\paragraph{A permission change not known at mount time.}
For two duties with no mounting release delay, both take $h$ to
mount and $c$ to move. One gets transfer permission at the
readout cut $\tau$; the other gets it at $\max(\tau,L)$.
Before $\tau$ the arbiter cannot distinguish which duty gets
which permission, cannot move either, and may still mount one.
The reliable readout and both times are public at the cut. The
arbiter may wait for it. Last move completion is the service
criterion, so a mount that sits in the center cannot pass.
A later update does not move or unmount an already placed duty.

If one duty is mounted before the cut, starting it at zero
weakly dominates any later precommitment. In the worst world
it needs the late permission; the other must wait for its full
physical departure. The worst completion time is
\[
 C=\max(h,\tau,L)+h+2c.
\]
If the arbiter has not mounted before the cut, start the early
duty then, using the order rule above:
\[
 W=\max(\tau+2h+c,L)+c.
\]
No third case can improve on these two. The worst case of
randomized policies is not smaller for a hard deadline, because
every supported action would have to meet that deadline in both
possible worlds. The exact bound is $\min(C,W)$, with waiting
strictly better precisely when $\tau+h<L$. At $\tau=0$ there is
no before-readout action; $W\le C$ and the minimum still applies.

For comparison, knowing which permission will arrive first
allows mounting that duty at zero, without taking its actual
permission early. Its optimum is
$\max(\max(h,\tau)+c+h,L)+c$. With $h=3,c=1,L=5,\tau=1$,
each world is feasible by tick 8 with that information.
An arbiter that must choose the first mount without it needs 10;
waiting lets both duties finish by tick 9. Thus two replacement
proofs that each pick a feasible early mount do not combine into
a lawful common control. Moving the readout here also changes the earliest usable
permission. The example does not isolate a pure information delay.

The direct tick-level search with the same information agrees
with the exact formulas. Lean checks the order exchange and
the strict waiting threshold. This is a conditional mechanism
for R2/L2 and a control for larger overlap tests; it is not a
claim of independent transfer, proof-work savings over a matched
scheduler, or contract closure.
\end{document}
